%! AP Statistics — Unit 1, Lesson 1.7 — Slides (Beamer)
\documentclass[aspectratio=169, 11pt]{beamer}
\usepackage{apstats-beamer}

%=============================================================================
\begin{document}

% ─────────────────────────────────────────────────────────────────────────────
% SLIDE 1 — LESSON TITLE
% ─────────────────────────────────────────────────────────────────────────────
{
\setbeamercolor{background canvas}{bg=navy}
\begin{frame}[plain]
  \begin{tikzpicture}[remember picture, overlay]
    \fill[navylight] (current page.north west)
      rectangle ([yshift=-1.35cm]current page.north east);
  \end{tikzpicture}
  \vspace{2.8cm}
  \hspace{0.55cm}%
  \begin{minipage}{0.9\textwidth}
    {\color{goldacc}\bfseries\small LESSON 1.7}\\[0.5cm]
    {\color{white}\bfseries\LARGE Summary Statistics for a\\[0.25cm]
      Quantitative Variable}\\[1.4cm]
    {\color{skymid}\small AP Statistics~$\cdot$~Unit 1: Exploring One-Variable Data}
  \end{minipage}
\end{frame}
}

% ─────────────────────────────────────────────────────────────────────────────
% SLIDE 2 — WARM-UP
% ─────────────────────────────────────────────────────────────────────────────
{
\setbeamercolor{background canvas}{bg=sky}
\begin{frame}[t, plain]
  \navyheader{Warm-Up}
  \vspace{1em}%
  \hspace{0.55cm}%
  \begin{minipage}{0.9\textwidth}
    \small
    \textit{Look at the dotplot on your warm-up sheet showing commute times for 10 students.}

    \vspace{0.5em}
    \begin{enumerate}
      \setlength{\itemsep}{0.6em}\setlength{\topsep}{0em}
      \item How many students are shown, and what is the variable?
      \item What is the range of commute times?
      \item List the values in order and identify the middle value(s).
      \item Which value stands apart from the rest? How can you tell from the dotplot?
    \end{enumerate}
    \vspace{0.4cm}
    {\color{navy}\itshape\small
      Today we give precise names to these observations: median, IQR, and the
      $1.5\times\text{IQR}$ outlier rule.}
  \end{minipage}
\end{frame}
}

% ─────────────────────────────────────────────────────────────────────────────
% SLIDE 3 — PRIMARY OBJECTIVE
% ─────────────────────────────────────────────────────────────────────────────
{
\setbeamercolor{background canvas}{bg=sky}
\begin{frame}[t, plain]
  \begin{tikzpicture}[remember picture, overlay]
    \fill[navy] (current page.north west)
      rectangle ([yshift=-0.25cm]current page.north east);
  \end{tikzpicture}
  \vspace{0.45cm}\par
  \hspace{0.55cm}%
  \begin{minipage}{0.9\textwidth}

    \sectionlabel{Primary Objective}

    \normalsize
    Students will \textbf{calculate and interpret} mean, median, range, IQR,
    and standard deviation, and will \textbf{choose the appropriate} measure of
    center and spread based on the shape of a distribution.

    \vspace{0.5cm}
    \textbf{Learning Targets:}
    \begin{enumerate}\small
      \setlength{\itemsep}{0.4em}
      \item Calculate mean and median; explain what each measures.
      \item Compute range, $Q_1$, $Q_3$, IQR, and standard deviation.
      \item Apply the $1.5\times\text{IQR}$ fence rule to identify outliers.
      \item Choose median + IQR vs.\ mean + $s_x$ based on shape and outliers.
    \end{enumerate}
  \end{minipage}
\end{frame}
}

% ─────────────────────────────────────────────────────────────────────────────
% SLIDE 4 — HOOK: THE SALARY PROBLEM
% ─────────────────────────────────────────────────────────────────────────────
{
\setbeamercolor{background canvas}{bg=goldbg}
\begin{frame}[t, plain]
  \navyheader{Hook: Which Salary Represents ``Typical''?}
  \vspace{1em}%
  \hspace{0.55cm}%
  \begin{minipage}{0.9\textwidth}
    \small
    Five friends report their annual salaries (thousands of dollars):
    \[\$30,\quad \$32,\quad \$35,\quad \$38,\quad \$200\]

    \vspace{0.4cm}
    \begin{columns}[t]
      \column{0.45\textwidth}
        \textbf{\color{navy}Mean:}
        \[\bar{x} = \frac{30+32+35+38+200}{5} = \$67\text{k}\]
        {\small The $\$200$k outlier \textit{pulls} the mean up.}

      \column{0.45\textwidth}
        \textbf{\color{navy}Median:}
        \[M = \$35\text{k}\]
        {\small The middle value is unaffected by $\$200$k.}
    \end{columns}

    \vspace{0.5cm}
    {\color{navy}\bfseries\small Discussion:} Which number would \textit{you} quote to
    describe a typical salary? Why does it matter which you choose?
  \end{minipage}
\end{frame}
}

% ─────────────────────────────────────────────────────────────────────────────
% SLIDE 5 — MEASURES OF CENTER
% ─────────────────────────────────────────────────────────────────────────────
{
\setbeamercolor{background canvas}{bg=white}
\begin{frame}[t, plain]
  \navyheader{Measures of Center}
  \vspace{1em}%
  \hspace{0.55cm}%
  \begin{minipage}{0.9\textwidth}
    \begin{columns}[t]
      \column{0.45\textwidth}
        \small
        \textbf{\color{navy}Mean} ($\bar{x}$):
        \[\bar{x} = \frac{\sum x_i}{n}\]
        \begin{itemize}\setlength{\itemsep}{0.4em}
          \item Uses \textit{every} data value
          \item \textbf{Not resistant} to outliers
          \item Pulled toward the \textit{tail} in skewed distributions
        \end{itemize}

      \column{0.45\textwidth}
        \small
        \textbf{\color{navy}Median} ($M$):
        \begin{itemize}\setlength{\itemsep}{0.4em}
          \item Sort data; take the middle value
          \item If $n$ even: average the two middle values
          \item \textbf{Resistant} to outliers
          \item Stays near the true center
        \end{itemize}
    \end{columns}

    \vspace{0.5cm}
    {\color{navy}\bfseries\small Key point:}
    When mean $>$ median $\Rightarrow$ skewed \textit{right}.\quad
    When mean $<$ median $\Rightarrow$ skewed \textit{left}.
  \end{minipage}
\end{frame}
}

% ─────────────────────────────────────────────────────────────────────────────
% SLIDE 6 — MEASURES OF SPREAD: RANGE & IQR
% ─────────────────────────────────────────────────────────────────────────────
{
\setbeamercolor{background canvas}{bg=white}
\begin{frame}[t, plain]
  \navyheader{Measures of Spread: Range \& IQR}
  \vspace{1em}%
  \hspace{0.55cm}%
  \begin{minipage}{0.9\textwidth}
    \begin{columns}[t]
      \column{0.45\textwidth}
        \small
        \textbf{\color{navy}Range:}
        \[\text{Range} = \max - \min\]
        \begin{itemize}\setlength{\itemsep}{0.3em}
          \item Simple; easy to compute
          \item \textbf{Not resistant} (uses only the extremes)
        \end{itemize}

        \vspace{0.4cm}
        \textbf{\color{navy}Quartiles:}
        \begin{itemize}\setlength{\itemsep}{0.3em}
          \item $Q_1$ = median of the \textit{lower} half
          \item $Q_3$ = median of the \textit{upper} half
          \item $\text{IQR} = Q_3 - Q_1$
        \end{itemize}

      \column{0.45\textwidth}
        \small
        \textbf{\color{navy}IQR captures the middle 50\%:}
        \begin{itemize}\setlength{\itemsep}{0.3em}
          \item \textbf{Resistant} to outliers
          \item Pair with \textit{median} when distribution is skewed
        \end{itemize}

        \vspace{0.4cm}
        \textbf{\color{navy}Example} (sorted): 3, 5, 7, 8, 9, 10, 12, 14, 16, 22\\[4pt]
        $Q_1 = 7$, $Q_3 = 14$, IQR $= 7$\\[4pt]
        Range $= 22 - 3 = 19$
    \end{columns}
  \end{minipage}
\end{frame}
}

% ─────────────────────────────────────────────────────────────────────────────
% SLIDE 7 — STANDARD DEVIATION
% ─────────────────────────────────────────────────────────────────────────────
{
\setbeamercolor{background canvas}{bg=white}
\begin{frame}[t, plain]
  \navyheader{Standard Deviation ($s_x$)}
  \vspace{1em}%
  \hspace{0.55cm}%
  \begin{minipage}{0.9\textwidth}
    \begin{columns}[t]
      \column{0.50\textwidth}
        \small
        \textbf{\color{navy}Formula:}
        \[s_x = \sqrt{\frac{\sum(x_i - \bar{x})^2}{n-1}}\]
        \begin{itemize}\setlength{\itemsep}{0.4em}
          \item Typical \textit{distance} from the mean
          \item Units match the original data
          \item Always $s_x \geq 0$; equals 0 only if all values identical
          \item Uses $n-1$ (sample standard deviation)
        \end{itemize}

      \column{0.42\textwidth}
        \small
        \textbf{\color{navy}Interpretation:}
        \begin{itemize}\setlength{\itemsep}{0.4em}
          \item Larger $s_x$ $=$ more spread out
          \item \textbf{Not resistant} to outliers
          \item Pair with \textit{mean} when distribution is symmetric
        \end{itemize}
        \vspace{0.3cm}
        \textbf{\color{navy}Remember:} Squaring the deviations makes them all positive;
        the square root returns units to the original scale.
    \end{columns}
  \end{minipage}
\end{frame}
}

% ─────────────────────────────────────────────────────────────────────────────
% SLIDE 8 — OUTLIER RULE
% ─────────────────────────────────────────────────────────────────────────────
{
\setbeamercolor{background canvas}{bg=sky}
\begin{frame}[t, plain]
  \navyheader{The $1.5 \times \text{IQR}$ Outlier Rule}
  \vspace{1em}%
  \hspace{0.55cm}%
  \begin{minipage}{0.9\textwidth}
    \small
    \begin{columns}[t]
      \column{0.50\textwidth}
        \textbf{\color{navy}Fence formulas:}
        \begin{align*}
          \text{Lower fence} &= Q_1 - 1.5(\text{IQR})\\
          \text{Upper fence} &= Q_3 + 1.5(\text{IQR})
        \end{align*}
        A value is an \textbf{outlier} if it falls\\
        \textit{below} the lower fence or\\
        \textit{above} the upper fence.

        \vspace{0.3cm}
        \textbf{\color{navy}Example:} $Q_1 = 7$, $Q_3 = 14$, IQR $= 7$\\
        Lower $= 7 - 10.5 = -3.5$\\
        Upper $= 14 + 10.5 = 24.5$\\
        Is 30 an outlier? \textbf{Yes} ($30 > 24.5$).

      \column{0.42\textwidth}
        \textbf{\color{navy}Important:}
        \begin{itemize}\setlength{\itemsep}{0.5em}
          \item An outlier flagged by this rule should be \textit{investigated}, not automatically removed.
          \item Check for data entry errors first.
          \item If real, report it --- explain its effect on the statistics.
        \end{itemize}
    \end{columns}
  \end{minipage}
\end{frame}
}

% ─────────────────────────────────────────────────────────────────────────────
% SLIDE 9 — CHOOSING THE RIGHT STATISTICS
% ─────────────────────────────────────────────────────────────────────────────
{
\setbeamercolor{background canvas}{bg=goldbg}
\begin{frame}[t, plain]
  \navyheader{Choosing the Right Statistics}
  \vspace{1em}%
  \hspace{0.55cm}%
  \begin{minipage}{0.9\textwidth}
    \small
    \begin{center}
    \renewcommand{\arraystretch}{1.6}
    \begin{tabular}{p{0.30\textwidth} p{0.28\textwidth} p{0.28\textwidth}}
    \rowcolor{navy}\color{white}\textbf{Situation} &
    \color{white}\textbf{Center} &
    \color{white}\textbf{Spread} \\
    \midrule
    Symmetric, no outliers & Mean ($\bar{x}$) & Std dev ($s_x$) \\
    Skewed \textit{or} has outliers & Median ($M$) & IQR \\
    \end{tabular}
    \end{center}

    \vspace{0.5cm}
    \begin{columns}[t]
      \column{0.45\textwidth}
        \textbf{\color{navy}Why pair them?}\\[4pt]
        Mean and standard deviation are both \textit{non-resistant} --- use them together
        when no outliers distort either.

      \column{0.45\textwidth}
        \textbf{\color{navy}Common AP error:}\\[4pt]
        Reporting mean + IQR, or median + $s_x$, is a mismatch.
        Always pair \textit{resistant} with \textit{resistant}.
    \end{columns}
  \end{minipage}
\end{frame}
}

% ─────────────────────────────────────────────────────────────────────────────
% SLIDE 10 — WORKED EXAMPLE (SWIM TIMES)
% ─────────────────────────────────────────────────────────────────────────────
{
\setbeamercolor{background canvas}{bg=white}
\begin{frame}[t, plain]
  \navyheader{Worked Example: Swim Team Lap Times}
  \vspace{1em}%
  \hspace{0.55cm}%
  \begin{minipage}{0.9\textwidth}
    \small
    Sorted lap times (sec), $n = 15$:\\
    62, 64, 65, 67, 68, 69, 70, \textbf{71}, 72, 73, 74, 75, 77, 80, 95

    \vspace{0.4cm}
    \begin{columns}[t]
      \column{0.45\textwidth}
        Median $= 71$ (8th value)\\[4pt]
        $Q_1 = 67$\quad $Q_3 = 75$\quad IQR $= 8$\\[4pt]
        Mean $= 1082 / 15 \approx 72.1$\\[4pt]
        Range $= 95 - 62 = 33$

      \column{0.45\textwidth}
        Fences:\\
        Lower $= 67 - 12 = 55$\\
        Upper $= 75 + 12 = 87$\\[4pt]
        Is 95 an outlier? \textbf{Yes} ($95 > 87$)\\[4pt]
        Mean $(72.1) >$ Median $(71)$ $\Rightarrow$ skewed right\\[4pt]
        $\Rightarrow$ Report \textbf{Median + IQR}
    \end{columns}
  \end{minipage}
\end{frame}
}

% ─────────────────────────────────────────────────────────────────────────────
% SLIDE 11 — GROUP ACTIVITY
% ─────────────────────────────────────────────────────────────────────────────
{
\setbeamercolor{background canvas}{bg=sky}
\begin{frame}[t, plain]
  \navyheader{Group Activity}
  \vspace{1em}%
  \hspace{0.55cm}%
  \begin{minipage}{0.9\textwidth}
    \small
    \begin{enumerate}
      \setlength{\itemsep}{0.6em}\setlength{\topsep}{0em}
      \item[\textbf{Part 1}] Calculate mean, median, $Q_1$, $Q_3$, IQR, and range
            for the swim team dataset.
      \item[\textbf{Part 2}] Apply the fence rule: is 95 an outlier? Which statistics
            should the coach report? Write a justification sentence.
      \item[\textbf{Part 3}] \textit{Individual:} A second coach reports mean $= 42$,
            median $= 44$. What shape does this suggest? Which center should be reported?
    \end{enumerate}

    \vspace{0.5cm}
    {\color{navy}\itshape\small
      Be ready to share: How did you decide which pair of statistics to use?}
  \end{minipage}
\end{frame}
}

% ─────────────────────────────────────────────────────────────────────────────
% SLIDE 12 — EXIT TICKET
% ─────────────────────────────────────────────────────────────────────────────
{
\setbeamercolor{background canvas}{bg=sky}
\begin{frame}[t, plain]
  \navyheader{Exit Ticket}
  \vspace{1em}%
  \hspace{0.55cm}%
  \begin{minipage}{0.9\textwidth}
    \small
    Hours of sleep reported by 8 students the night before an exam (sorted):
    \[4,\quad 7,\quad 7,\quad 9,\quad 10,\quad 13,\quad 15,\quad 30\]

    \vspace{0.4em}
    \begin{enumerate}
      \setlength{\itemsep}{0.5em}\setlength{\topsep}{0em}
      \item Find the mean and median. Show calculations.
      \item Find $Q_1$, $Q_3$, and IQR.
      \item Use the fence rule: is 30 an outlier?
      \item \textbf{AP-Style:} Which measure of center is more appropriate?
            Write one complete sentence with justification and context.
    \end{enumerate}

    \vspace{0.4cm}
    {\color{navy}\itshape\small
      Complete individually --- 5 minutes.}
  \end{minipage}
\end{frame}
}

% ─────────────────────────────────────────────────────────────────────────────
% SLIDE 13 — WRAP-UP & PREVIEW
% ─────────────────────────────────────────────────────────────────────────────
{
\setbeamercolor{background canvas}{bg=navy}
\begin{frame}[plain]
  \vspace{1.6cm}
  \hspace{0.55cm}%
  \begin{minipage}{0.9\textwidth}
    {\color{goldacc}\bfseries\large Lesson 1.7 Summary}\\[0.6cm]
    \small\color{white}
    \begin{itemize}\setlength{\itemsep}{0.5em}
      \item \textbf{Mean} uses all values; pulled by outliers.
            \textbf{Median} uses position; resistant.
      \item \textbf{IQR} $= Q_3 - Q_1$; resistant measure of middle 50\%.
      \item \textbf{Outlier rule:} below $Q_1 - 1.5\cdot\text{IQR}$ or above
            $Q_3 + 1.5\cdot\text{IQR}$.
      \item \textbf{Decision:} Skewed/outliers $\to$ Median + IQR;\quad
            Symmetric $\to$ Mean + $s_x$.
    \end{itemize}
    \vspace{0.6cm}
    {\color{skymid}\bfseries\small Preview:}
    {\color{white}\small Lesson 1.8 organizes these five numbers
    ($\min,\,Q_1,\,M,\,Q_3,\,\max$) into a \textbf{boxplot} --- a powerful
    visual for comparing distributions.}
  \end{minipage}
\end{frame}
}

\end{document}
